% !PRINTED 663
\weberpage{663}
%<<CONT
also nach (8)
\weq{10}{\Phi_1^n \Phi_n^{-1} = \prod^{t} \varphi_{t'}^{\,m E\left(\frac{n t}{m}\right)} = \varepsilon\, \alpha^m.}

Hiernach ist der Quotient
\[
\left(\frac{\prod\limits^{t} \varphi_{t'}^{\,E\left(\frac{n t}{m}\right)}}{\alpha}\right)^m
\]
ein Einheitsfunctional, und da er zugleich die $m^{\text{te}}$ Potenz eines
Functionals in $\Omega_m$ ist, so ist die Basis dieser Potenz auch eine
Einheit, d.\,h. es ist
\[
\prod^{t} \varphi_{t'}^{\,E\left(\frac{n t}{m}\right)} = \varepsilon \alpha,
\]
und es gehört dies Product der Hauptclasse an. Nun ist aber
nach unserer Voraussetzung
\[
\varphi_{t'}\ \text{äquivalent mit}\ \varphi^{t'},
\]
und daher ist
\[
\prod^{t} \varphi_{t'}^{\,E\left(\frac{n t}{m}\right)}\ \text{äquivalent mit}\ \varphi^{\sum^{t} t' E\left(\frac{n t}{m}\right)},
\]
und dies gilt für jedes beliebige durch $q$ nicht theilbare $n$.

Können wir also nachweisen, dass sich $n$ so bestimmen
lässt, dass auch die Summe
\[
S_n = \sum^{t} t' E\left(\frac{n t}{m}\right)
\]
durch $q$ nicht theilbar ist, so haben wir das Theorem a) bewiesen.

Wir wollen zunächst eine Umformung der Summe $S_n$ vornehmen,
durch die die Frage auf den weit einfacheren Fall
zurückgeführt wird, in dem $m$ eine Primzahl ist.

Wir setzen, wenn $m$ eine höhere als die erste Potenz von
$q$ ist,
\[
m = q m', \quad t = t_1 + q t_2, \quad
\begin{aligned}
t_1 &= 1, 2, \dots, q-1 \\
t_2 &= 0, 1, \dots, m'-1.
\end{aligned}
\]

Darin ist $m'$ noch eine Potenz von $q$, der Summationsbuchstabe
 $t_1$ durchläuft die Zahlenreihe $1, 2, \dots, q-1$, und $t_2$
die Zahlenreihe $0, 1, 2, \dots, m'-1$. Wir bestimmen $t_1'$ aus der
Congruenz
\[
t_1 t_1' \equiv 1\ (\mathrm{mod}\ q)
\]
und erhalten
\[
t' \equiv t_1'\ (\mathrm{mod}\ q).
\]
